%% ma: L12-B14-0114
%% lop: 12
%% chuong: 5
%% bai: 14
%% dang: 12.14.32
%% loai: TN4
%% muc: TH
%% dap_an: "A"
%% nguon: "Ảnh giáo viên gửi 07/11/2026 (Tư Duy Mở - tuduymo.com, đề tự luyện Phương trình mặt phẳng, đề khoảng cách), Câu 10"
%% kien_thuc: "Bài 14 (hai mặt phẳng song song, khoảng cách từ một điểm đến mặt phẳng)"
%% trang_thai: cho-duyet
%% ly_do: "Dạng 12.14.32 mới đề xuất, chờ giáo viên duyệt."
%% ghi_chu: "Tài liệu gốc không kèm đáp án; đáp án do Claude giải và kiểm bằng sympy."
\begin{ex}
Trong không gian $Oxyz$, cho mặt phẳng $(P)\colon x-2y+2z-1=0$ và điểm $A(1;0;-1)$. Có hai mặt phẳng phân biệt $(Q_1)$ và $(Q_2)$ song song với $(P)$ và cách $A$ một đoạn bằng $3$. Hai mặt phẳng $(Q_1)$ và $(Q_2)$ cách nhau một khoảng là
\choice{\True $6$}{$3$}{$\dfrac32$}{$0$}
\loigiai{$(Q_1)$, $(Q_2)$ song song và nằm về hai phía của $A$, mỗi mặt cách $A$ một khoảng $3$, nên cách nhau $6$. (Cụ thể $x-2y+2z+10=0$ và $x-2y+2z-8=0$.) Chọn A.}
\end{ex}
